% ============================================================
%  Capítulo 5 -- Planificación y Gestión
% ============================================================

\section{Planificación y Gestión}

El desarrollo de OmniLitter se ha planteado de forma iterativa, combinando tareas
de análisis, diseño, prototipado e implementación. Esta aproximación resulta
adecuada para un proyecto con cliente real, ya que permite validar decisiones antes
de invertir esfuerzo en una implementación definitiva.

\subsection{Metodología de trabajo}

La metodología seguida se basa en ciclos cortos de trabajo:

\begin{enumerate}
  \item Recogida de necesidades con el cliente y tutores.
  \item Definición de requisitos funcionales, no funcionales y de datos científicos.
  \item Elaboración de mockups funcionales para validar flujos de usuario.
  \item Revisión con el cliente y refinamiento del prototipo.
  \item Diseño de la arquitectura, base de datos y componentes del sistema.
  \item Implementación progresiva de las partes finales del proyecto.
\end{enumerate}

Aunque no se ha aplicado una metodología ágil formal completa, el proyecto adopta
una forma de trabajo incremental: cada avance genera un artefacto revisable
(documento de requisitos, mockup, diagrama, módulo de código o capítulo de memoria)
que se contrasta antes de continuar.

\subsection{Gestión de reuniones}

Las reuniones con el cliente han tenido un papel central en la definición del
alcance. La primera reunión permitió identificar el problema y las funcionalidades
esperadas. Posteriormente, una segunda reunión sirvió para revisar los mockups,
confirmar la cobertura de los flujos principales y reforzar prioridades como el
funcionamiento sin conexión, la separación por grupos y la trazabilidad científica.

Estas reuniones se documentan en el Capítulo~6 y en el Anexo~\ref{sec:anexo_requisitos_detallados},
donde se relacionan las necesidades detectadas con los requisitos del sistema.

\subsection{Gestión del alcance}

El alcance se ha dividido en dos niveles:

\begin{itemize}
  \item \textbf{Alcance funcional validado mediante mockups:} portal web y aplicación
    móvil con flujos navegables, datos de ejemplo, roles simulados y pantallas
    principales implementadas.
  \item \textbf{Alcance de implementación final:} backend REST, base de datos,
    autenticación, sincronización offline, persistencia de observaciones y conexión
    progresiva de los clientes al backend.
\end{itemize}

Esta separación permite demostrar primero la experiencia de uso y después construir
la infraestructura técnica necesaria para convertir el prototipo en una aplicación
real.
